% !TeX root = ./main.tex
\section{Methodology}
\label{sec:methods}

\subsection{Problem formulation}
\label{sec:problem}

The forecasting task is posed as sequence-to-sequence regression on the two-dimensional (frequency $\times$ time) directional wave spectrum measured at a single buoy. At each hourly time step $t$, the buoy record consists of
five co-located quantities defined over a fixed frequency grid $f_1,\dots,f_F$ ($F=47$ bins spanning $0.0200$--$0.4850$~Hz, non-uniformly spaced: finer resolution at low frequencies, coarser above $\sim$0.3~Hz): the one-dimensional spectral density $E(f,t)$ [m$^2$/Hz], the mean wave direction $\alpha_1(f,t)$, the principal wave direction $\alpha_2(f,t)$, and the first and second normalised polar coordinates (directional spreading parameters) $r_1(f,t)$ and $r_2(f,t)$. The pair $r_1$, $r_2$ jointly bounds the directional spreading function, and a ratio $r_2/r_1^2$ departing from unity flags mixed or bimodal sea states that $r_1$ alone cannot distinguish; $r_2$ is therefore retained as an input channel rather than discarded as redundant.

Because $\alpha_1$ and $\alpha_2$ are circular quantities (degrees), each is presented to the model as a $(\sin,\cos)$ pair rather than as a raw angle, so that e.g.\ $1^\circ$ and $359^\circ$ are adjacent in input space instead of maximally separated. The encoder therefore observes $C=7$ channels per frequency bin,
$\mathbf{x}(f_k,t) = [\,E,\ \sin\alpha_1,\ \cos\alpha_1,\ \sin\alpha_2,\ \cos\alpha_2,\ r_1,\ r_2\,](f_k,t)$.

Given an input window of the $L$ most recent hourly observations of all seven channels, $\mathbf{X}_t = \mathbf{x}(f_{1:F},\, t-L+1{:}t)$, the model forecasts a spectral trajectory over a forecast horizon of $\tau$ hourly steps, for $\tau \in \{12, 24, 48\}$~h. Each horizon is treated as an independent forecasting problem: one model per horizon is trained and hyperparameter-tuned separately (Section~\ref{sec:training}), each producing its own full $\tau$-step-ahead trajectory rather than a single terminal-step prediction, although — unlike in the earlier configuration of this work — only the terminal step is treated as the forecast product, the intermediate steps being autoregressive scaffolding required to reach it (Section~\ref{sec:evaluation}).

\paragraph{Shape/magnitude decomposition of the forecast target}
Rather than forecasting the physical spectrum $E(f,t)$ directly, the forecast target is factorised into a shape and a magnitude component, which are modelled separately and recombined at inference time:
\begin{equation}
E(f,t) \;=\; S(f,t)\, m_0(t), \qquad
S(f,t) = \frac{E(f,t)}{m_0(t)}, \qquad
m_0(t) = \Big(\tfrac{H_s(t)}{4}\Big)^{2},
\label{eq:shape_magnitude}
\end{equation}
where $S$ is the unit-area spectrum ($\int S(f)\,\mathrm{d}f = 1$ by construction) and $m_n = \int f^n E(f)\,\mathrm{d}f$ are the spectral moments. A \texttt{shape}-target model forecasts $S(f,\,t+1{:}t+\tau)$ and a companion \texttt{hs}-target model forecasts the scalar $H_s = 4\sqrt{m_0}$; their outputs are multiplied per Eq.~\eqref{eq:shape_magnitude} to reconstruct the physical spectrum. This decouples the (comparatively easy) magnitude-forecasting problem from the (harder) shape-forecasting problem, and is motivated by the tendency of a single monolithic $E(f)$-target model to underestimate spectral peaks and over-smooth the high-frequency tail. A monolithic \texttt{density} target ($E(f)$ forecast directly) remains implemented and is used as a reference configuration for the comparison in Section~[INSERT: Results section reference]; unless stated otherwise, results below refer to the shape/magnitude split. Because the two branches share the same data pipeline, architecture, training procedure and evaluation protocol — differing only in output dimensionality ($F$ versus $1$) and in the target transform described in Section~\ref{sec:data} — the description that follows is written once and applies to both.

Both spectral targets are predicted in logarithmic space: the network's output layer emits $\log S(f)$ (or $\log E(f)$ for the \texttt{density} configuration) rather than the quantity itself, so that non-negativity of the recovered physical spectrum is obtained for free by exponentiation at the point of use, with no architectural constraint (such as a Softplus output activation) and no clipping. Bulk sea-state parameters, in particular $H_s$ and the mean zero-crossing period $T_{m02} = \sqrt{m_0/m_2}$, are derived from the forecast spectrum post hoc and used for evaluation (Section~\ref{sec:evaluation}); $H_s$ is additionally a training target in its own right for the magnitude branch, per Eq.~\eqref{eq:shape_magnitude}.

\subsection{Data}
\label{sec:data}

\paragraph{Source}
Data were obtained from a single directional wave buoy (NDBC station 32012; Woods Hole Stratus Oceanographic Institution, (19°25'29" S 85°4'39" W)), recording the five channels described in Section~\ref{sec:problem} at hourly resolution. The record used here spans 2016-01-01 00:00 to 2017-12-31 23:00 (17,544 hourly observations, $\approx$2 years).

\paragraph{Preprocessing}
All five channels are reindexed onto a single shared hourly grid spanning the union of their timestamps, and any resulting gaps are filled by linear-in-time interpolation (bidirectional, so leading and trailing gaps are also filled). The two directional channels $\alpha_1$ and $\alpha_2$ are interpolated on the unit circle rather than on the raw angle: each is decomposed into $\sin$/$\cos$ components per frequency bin, the components are interpolated in time independently, and the angle is recovered with $\operatorname{atan2}$. Interpolating the raw angle would traverse the wrong side of the compass across a gap (e.g.\ $350^\circ \to 10^\circ$ passing through $180^\circ$ instead of through $0^\circ$). Because gaps are interpolated away at this stage, the model is never exposed to irregular time sampling: positional encoding (Section~\ref{sec:architecture}) therefore encodes a fixed, regularly-spaced sequence index rather than solving for irregularly-sampled inputs.

\paragraph{Train/validation/test split}
The split is strictly temporal and chronological, with no shuffling across split boundaries: the first 70\% of the hourly record is used for training (2016-01-01 00:00 -- 2017-05-26 15:00), the next 15\% for
validation (2017-05-26 16:00 -- 2017-09-13 07:00), and the final 15\% for test (2017-09-13 08:00 -- 2017-12-31 23:00). The test split is used only once, after hyperparameter search, to report final model performance (Section~\ref{sec:training}). This is a temporal hold-out only. No spatial (unseen-buoy) generalisation is assessed.

\paragraph{Normalisation}
All normalisation statistics are fit on the training split only and applied unchanged to validation and test. The spectral density channel $E(f,t)$ only uses scale normalisation: $\tilde{E}(f,t) = E(f,t)/\mu(f)$, where $\mu(f)$ is
the per-frequency training-set mean (floored at $10^{-8}$ to avoid division by zero). This choice preserves non-negativity of the spectrum, which is required because $H_s$ is computed as $4\sqrt{m_0}$ from the (denormalised) spectrum. This normalise/denormalise round trip, and its effect on $H_s$, is verified directly in the test suite (JONSWAP-spectrum round-trip test, Section~\ref{sec:evaluation}). The directional channels carry no non-negativity constraint and are handled according to their own geometry: $r_1$ and $r_2$ are $z$-scored (training mean subtracted, divided by training standard deviation, floored at $10^{-8}$), whereas the $\sin$/$\cos$ decompositions of $\alpha_1$ and $\alpha_2$ are passed through unnormalised, since they already lie in $[-1,1]$.

\paragraph{Target construction and log transform}
Targets are built from the physical (pre-normalisation) density, so that the target series and the persistence baseline are meaningful in physical units without reference to $\mu(f)$: the target is the unit-area spectrum $S(f) = E(f)/m_0$ (with $m_0$ floored at $10^{-12}$~m$^2$). Both spectral targets are then mapped into log space, $y = \log \max\!\big(y_{\text{phys}},\, \epsilon\,\bar{y}(f)\big)$, where $\bar{y}(f)$ is the per-frequency training-split mean of the corresponding physical target and $\epsilon = 10^{-3}$. The per-bin floor, rather than a single global constant, keeps the floor proportional to the energy each bin actually carries, and prevents $\log$ from being evaluated at exactly zero for a genuinely calm bin. The same transform is applied to the decoder's start token and to every teacher-forced or self-generated decoder input, so that the entire autoregressive loop operates consistently in log space.

\paragraph{Windowing}
Input/target pairs are built by sliding a window of $L$ consecutive hourly steps (input) followed by $\tau$ consecutive hourly steps (target) across each split, yielding $T_{\text{split}} - L - \tau + 1$ samples per split, where $T_{\text{split}}$ is the number of hourly records in that split. The input window length $L$ is itself a tuned hyperparameter (Section~\ref{sec:training}) rather than a fixed constant, so realised sample counts vary by trial.

\paragraph{Auxiliary dynamical features}
In addition to the frequency-resolved channels above, each sample carries a small set of scalar side features summarising the dynamics of its own input window, obtained by Dynamic Mode Decomposition (DMD)~\cite{schmid2010dmd}. For each sample independently, a linear operator $A$ satisfying $\mathbf{x}_{t+1} \approx A\,\mathbf{x}_t$ is fitted across the $L$ snapshots of that window's (normalised) density spectra, and $A$ is eigendecomposed; each eigenvalue yields a growth/decay rate (from its magnitude) and an oscillation frequency (from its phase), together with the corresponding mode amplitude. The four dominant modes by amplitude are retained (giving $4 \times 3 = 12$ scalar features per sample) $z$-scored on the training split and broadcast across the input window. The motivation is to state explicitly, rather than requiring the attention layers to infer it implicitly, whether the currently observed wave systems are growing or decaying. These features are not frequency-resolved and are fused into the encoder through a separate pathway (Section~\ref{sec:architecture}); they are encoder-only, since they are never a forecast target. DMD eigenvalues are invariant to a fixed per-bin rescaling of the state (a similarity transform), so computing them on the normalised rather than the physical spectrum is not a loss of information.

\subsection{Model architecture}
\label{sec:architecture}

The model is a standard encoder--decoder Transformer~\cite{vaswani2017attention} with multi-head scaled dot-product attention, applied to the multi-channel spectral sequence described above, in its pre-normalisation (\texttt{norm\_first}) variant, which is more stable to train at the deeper end of the search space than the original post-norm arrangement.

\paragraph{Encoder input embedding}
Rather than flattening the $(F \times C)$ per-timestep measurement through a single unconstrained linear layer, the encoder embedding respects the (frequency, channel) structure of the input and maps it to one temporal token in four stages:
\begin{enumerate}
  \item \emph{Per-bin projection.} A linear layer shared across all $F$ frequency bins projects the $C=7$ channel measurement at each bin into an 8-dimensional per-bin representation (GELU-activated). The rationale is physical: the channels at a given frequency bin (energy and the directional parameters) describe the same underlying process at that frequency and are therefore combined before being mixed across frequencies, rather than being treated as $F \times C$ independent scalars.
  \item \emph{Frequency-identity signal.} A fixed sinusoidal encoding of each bin's actual \emph{log-frequency} value is added to its representation, together with a learned per-bin residual initialised at zero (so training begins from the purely physical signal and only learns corrections the grid alone does not capture). The sinusoid's wavelengths are scaled to the observed span of $\log f$ so that distance in embedding space tracks true log-frequency distance rather than array position. This matters because the grid is non-uniform: a fixed index offset corresponds to $\sim$0.005~Hz near 0.02~Hz and to $\sim$0.02~Hz near 0.485~Hz.
  \item 
  \item \emph{Frequency-axis convolution.} The per-bin representations pass through the same dilated convolutional block used on the time axis (described below), applied along the frequency axis, so each bin absorbs local context from its neighbours — exploiting the fact that a wave spectrum is a smooth curve in frequency. Padding is replicative rather than zero here, because the frequency axis is bounded and non-cyclic: zero padding fabricates fictitious zero-energy bins immediately outside the $0.02$--$0.485$~Hz grid and measurably corrupts the bins nearest each edge.
  \item \emph{Attention pooling.} A single learned query attends over the frequency axis and collapses all $F$ bins into one $d_{\text{model}}$-dimensional token. Unlike a flatten-and-project aggregation, the softmax weighting is content-dependent, as it can follow wherever the spectral peak currently sits, instead of applying a per-bin-position weight fixed at training time — and its parameter count is independent of $F$.
\end{enumerate}
The empirical comparison against the flat-projection alternative is reported in Section~[INSERT: Results section reference].

\paragraph{Auxiliary feature fusion}
The DMD features of Section~\ref{sec:data} are not frequency-resolved and therefore bypass the embedding above entirely: they are projected by their own linear layer to $d_{\text{model}}$ and added into the per-timestep encoder token, before positional encoding. This is the same additive convention positional encoding itself uses. They are supplied to the encoder only, never to the decoder.

\paragraph{Positional encoding}
A fixed sinusoidal positional encoding~\cite{vaswani2017attention} is added to both encoder and decoder token sequences (shared weights, i.e. the same encoding table). Because attention has no inherent notion of sequence order and the input sequence is, by construction (Section~\ref{sec:data}), strictly regular in time, a fixed encoding of the discrete step index is sufficient; no mechanism for irregular or missing-timestamp inputs is required or implemented.

\paragraph{Temporal convolutional front-end (encoder only)}
Following positional encoding, the encoder sequence passes through three dilated 1-D causal-style convolutions (kernel size 3; dilations 1, 2, 4; padding equal to the dilation, preserving sequence length), each wrapped in a pre-norm residual block with GELU activation and dropout. The combined stack has an effective receptive field of 9 consecutive hourly steps. This front-end injects local temporal structure (e.g. a swell rising over a few hours) into each token before the global self-attention layers, so that attention capacity is spent on longer-range dependencies rather than re-discovering short, local trends. It is applied only to the encoder sequence: the decoder sequence length equals the forecast horizon $\tau$ ($\leq$48 steps) and already carries a causal mask, so the same local-context argument does not apply.

\paragraph{Decoder}
The decoder receives a single-channel spectrum per step and embeds it through the \emph{same} four-stage frequency-structured embedding as the encoder, instantiated with $C=1$, so that both sides of the model are given the same structural prior over the frequency axis. The decoder is seeded with a start token equal to the last (most recent) spectrum of the input window, denormalised, converted to the target quantity, and log-transformed exactly as the targets are; at inference, subsequent tokens are the model's own previous-step predictions (autoregressive rollout), while a subsequent-position (causal) mask restricts each decoder position to attend only to earlier positions during training. The final linear layer projects each decoder position back to $\mathbb{R}^{F}$, producing one forecast log-spectrum per forecast step. Each predicted step is exponentiated and renormalised to unit area (by trapezoidal integration over the physical frequency grid) before being fed back as the next decoder input, so that the autoregressive state remains a valid probability density throughout the rollout; the encoder is evaluated once per sample and its output reused across all $\tau$ decode steps.

\paragraph{Capacity}
The model dimension is constructed as $d_{\text{model}} = d_{\text{head}}
\times n_{\text{head}}$, with $d_{\text{head}}$ and $n_{\text{head}}$ fixed at $d_{\text{head}}=32$, $n_{\text{head}}=8$ (hence $d_{\text{model}}=256$) rather than searched, on the basis of a tally across all previous searches of the same space, in which those two values won at a majority of (study, horizon) points while every other hyperparameter spanned nearly its entire range with no consensus; this frees two dimensions of the search budget for the loss weights 

\subsection{Training procedure}
\label{sec:training}

\paragraph{Loss function}


The spectral loss is a composite of three physically motivated terms evaluated on the log-space prediction $\hat{y}$ and target $y$,
\begin{equation}
\mathcal{L} \;=\; \lambda_{\text{base}}\,\mathcal{L}_{\text{bin}}
\;+\; \lambda_{\text{KL}}\,D_{\text{KL}}
\;+\; \lambda_{W}\,W_2
\;+\; \lambda_{P}\,\mathcal{L}_{\text{peak}} .
\label{eq:loss}
\end{equation}
The four terms are:
\begin{itemize}
  \item $\mathcal{L}_{\text{bin}}$ — a per-bin mean squared error between $\hat{y}$ and $y$ in log space, weighted across the frequency axis by normalised trapezoidal weights $w_i \propto \Delta f_i$. The weighting is necessary because the grid is log-spaced: a flat elementwise mean over bins over-represents the dense low-frequency region relative to its actual share of the physical integral.
  \item $D_{\text{KL}}$ — the Kullback--Leibler divergence~\cite{kullback1951} $D_{\text{KL}}(P_{\text{true}}\,\|\,P_{\text{pred}})$ between the two spectra treated as discrete probability distributions over frequency, with $P(f_i) \propto E(f_i)\,\Delta f_i$ ($\Delta f$-weighted for the same reason as above). Since both tensors are already logarithmic, both distributions are obtained directly as $\log P = \operatorname{log\_softmax}(y + \log \Delta f)$, which is unconditionally numerically stable. Because the divergence is weighted by $P_{\text{true}}$, it is dominated by the model's assigned mass \emph{at} the bins where the true spectrum is concentrated, and is comparatively indifferent to how the remaining mass is distributed among low-energy bins; a peak that broadens into its immediate neighbourhood while still covering its own bin is thus penalised far less than under a per-bin MSE, which charges separately for every newly elevated neighbour. KL is used in place of the gradient-equivalent cross-entropy purely so that the reported value is zero at a perfect match.
  \item $W_2$ — the 1-D Wasserstein-2 (quadratic earth-mover) distance between the two spectra, each first normalised by its own total mass so that the term measures shape alone, decoupled from magnitude. For 1-D distributions this is computed exactly in the quantile domain, $W_p^p = \int_0^1 |F^{-1}(q) - G^{-1}(q)|^p\,\mathrm{d}q$, with the true spectrum's own cumulative levels serving as quadrature nodes (the true quantile function inverts trivially at those levels, so only the predicted CDF requires inversion); no optimal-transport solver or critic network is needed~\cite{peyre2019}. A quadratic rather than linear transport cost is used so that one badly displaced peak is penalised disproportionately more than several small local shifts of the same total distance. Unlike a pointwise loss — which penalises a peak displaced by one bin almost as harshly as one displaced entirely, since both give near-zero overlap at the peak location — $W_2$ is tolerant of small position errors while still charging real transport cost for replacing two sharp peaks with one flat blur.
  \item $\mathcal{L}_{\text{peak}}$ — a differentiable soft peak-height term. Within each trough-to-trough window $k$ of the \emph{true} spectrum (detected as described in Section~\ref{sec:evaluation}), the predicted peak height is estimated as a temperature-weighted softmax average of the predicted spectrum against itself, $\hat{H}_k = \sum_{i \in \text{window}_k} \hat{E}(f_i)\,\sigma_k(i;\tau_k)$ with $\sigma_k(i;\tau_k) \propto \exp(\hat{E}(f_i)/\tau_k)$, and the loss is the mean over peaks of $(\hat{H}_k - H_k^{\text{true}})^2$, where $H_k^{\text{true}}$ is the exact maximum of the true spectrum within that window. The softmax replaces a hard $\max$ on the prediction side, whose subgradient is zero everywhere except at a winning bin that can flip discontinuously between steps; the label side needs no gradient and so keeps the exact maximum. As the soft estimate is bounded above by the hard maximum of the same values, this asymmetry leaves a small positive residual even at a perfect prediction (shrinking with $\tau_k$), which is accepted as a floor on this term rather than removed by also softening the label. The temperature is set per peak as $\tau_k = \max\!\big(\tau_{\min},\, c\,H_k^{\text{true}}\,(f_{r_k}-f_{l_k})/(f_F - f_1)\big)$ with $c=0.15$: it is scaled by the peak's own energy, because $\exp(E/\tau)$ is only meaningful when $\tau$ carries $E$'s units, and modulated by the window's bandwidth in Hz (not in bins, again because the grid is non-uniform), so that narrow swell partitions get a sharp temperature and broad wind-sea partitions a softer one. The term is invariant to pure translation of a peak and therefore complements $W_2$, which measures displacement but not amplitude loss.
\end{itemize}
The weights $\lambda_{\text{KL}}$, $\lambda_{W}$ and $\lambda_{P}$ are tuned per horizon (Section~\ref{sec:hpo}). The composition itself, however, is fixed rather than searched: $\lambda_{\text{base}} = 0$, i.e. the per-bin term is \emph{substituted} by the three distributional terms rather than augmented with them. This follows a dedicated, fixed-architecture ablation over seven candidate recipes, described in Section~\ref{sec:ablation}, in which the three-way combination was the only one that improved on the per-bin baseline across the whole evaluation panel rather than trading one axis for another.

Note that for the spectral branches the loss is evaluated in log space, on the quantity the network emits directly, so no denormalisation step enters the gradient path. This differs from the earlier physical-space formulation of this work, in which prediction and target were denormalised ($E = \tilde{E}\cdot\mu(f)$) before an RMSE was taken; the frequency weighting described above now plays the role that denormalisation previously played, i.e.\ preventing bins with small $\mu(f)$ from dominating the gradient irrespective of their physical energy content.

\paragraph{Optimisation}
Parameters are optimised with AdamW~\cite{loshchilov2019adamw}, which decouples weight decay from the adaptive gradient scaling rather than folding it in as an $L^2$ penalty; learning rate ($10^{-3}$--$1.5\times10^{-2}$, log-uniform) and weight decay ($10^{-6}$--$10^{-2}$, log-uniform) are tuned per Section~\ref{sec:hpo} rather than fixed. The learning rate is ramped linearly from zero to its sampled value over the first 5 epochs. Without this warmup, training at the upper end of the learning-rate range is unstable from the first epoch, which was diagnosed as the cause of an anomaly in earlier searches, in which the learning rate dominated hyperparameter-importance rankings while correlating only weakly with the achieved score — the signature of noisy early training rather than of a clean optimum the sampler can exploit. Gradients are clipped to a maximum norm of 1.0 at every step, to prevent early-training loss spikes from displacing parameters into poorly-conditioned regions of the loss surface. After warmup, the learning rate is reduced on plateau (factor 0.5) after 5 epochs without improvement in the validation objective (Section~\ref{sec:evaluation}), with a 2-epoch cooldown before the scheduler resumes counting, so that at least one reduction occurs well before early stopping can trigger (see below), giving optimisation a chance to benefit from the smaller step size before training halts. The scheduler is suppressed during warmup, so that the ramp itself is never mistaken for a plateau.

\paragraph{Scheduled sampling}
Training uses teacher forcing with a linearly decaying teacher-forcing ratio $p$, from $p=1$ at epoch 0 to $p=0$ at epoch $2\times\text{patience}$ (i.e. 40 epochs, with patience defined below): at each decoder step, the ground-truth previous token is fed with probability $p$ and the model's own (detached) prediction with probability $1-p$. This is the standard mitigation for exposure bias in autoregressive sequence models~\cite{bengio2015scheduledsampling}.  Without it, the decoder is trained exclusively on ground-truth history but evaluated autoregressively on its own (imperfect) predictions. The draw is made per sample rather than once per batch, so that every batch sees a mixture close to the target ratio; a single batch-wide draw would instead make whole batches uniformly teacher-forced or uniformly autoregressive by chance, injecting spurious variance into the training signal. The decay horizon is tied to the early-stopping patience (rather than a fixed constant) so that runs which are early-stopped later still see a materially reduced teacher-forcing ratio.

\paragraph{Early stopping}
Training runs for up to 100 epochs with early stopping at patience $=20$ epochs without improvement in the validation objective; the model weights from the best validation epoch (not the last epoch) are retained for
evaluation. Because the validation objective is itself computed by autoregressive rollout on a comparatively small split, its epoch-to-epoch value is noisy; it is therefore smoothed by a 5-epoch trailing mean before it drives any decision — learning-rate scheduling, checkpoint selection, early stopping, and (during the search) pruning alike. Using raw per-epoch values was found to lock checkpoint selection onto an early noise spike and to stop training on a stale, undertrained checkpoint while the underlying metric was still improving.

\subsubsection{Hyperparameter search}
\label{sec:hpo}

Hyperparameters are tuned with Optuna's~\cite{akiba2019optuna} Tree-structured Parzen Estimator sampler~\cite{bergstra2011tpe} (multivariate mode, 18 startup trials before the multivariate surrogate takes over, fixed seed). Pruning is performed by a median pruner that becomes active only after 30 of the maximum 100 epochs and then checks every 5 epochs (matching the smoothing window above, so that successive checks compare largely non-overlapping windows rather than re-examining the same slow-moving average); the 30-epoch warmup replaced an earlier 20-epoch one after it was found to prune a majority of trials at exactly the boundary epoch. The median verdict is further gated by a patience rule requiring 10 consecutive non-improving reports before a trial is actually terminated, because trials were observed to dip and recover several times — including the eventual best trial — around the warmup and teacher-forcing-decay boundaries.

The search space is: input window length $L \in \{12, 24, 48, 96\}$~h; batch size $\in \{32, 64, 128\}$, restricted to $\{32, 64\}$ for $\tau > 24$~h because the scheduled-sampling rollout's memory footprint grows quadratically with the horizon; learning rate and weight decay as above; dropout on the 8-dimensional per-bin representation $\in [0.1, 0.3]$ and dropout on the $d_{\text{model}}$-wide representation (positional encoding, temporal convolutions, and the Transformer's own attention/feed-forward dropout) $\in [0.0, 0.3]$, tuned separately because dropping units from the narrow per-bin representation is a much larger relative perturbation than dropping them from the wider token representation; number of encoder layers and decoder layers, each $\in \{1,2,3,4\}$; and the three loss weights of Eq.~\eqref{eq:loss}, $\lambda_{\text{KL}} \in [0.1, 50]$, $\lambda_{W} \in [1, 200]$ and $\lambda_{P} \in [0.01, 20]$, all log-uniform over the ranges validated by the loss ablation described above. Attention head dimension and head count are pinned as described in Section~\ref{sec:architecture}. Eighty trials are run per forecast horizon $\tau$, each trained and early-stopped independently as described above, as a separate job. The optimisation objective (Section~\ref{sec:evaluation}) is evaluated on the validation split only; the held-out test split is evaluated once per trial
purely for monitoring/logging and does not influence trial selection, which is driven solely by the validation-split objective.

The global random seed is fixed once, before the Optuna loop begins, and is deliberately \emph{not} reset inside each trial: resetting it per trial would force every trial to start from an identical initialisation and data order, collapsing the outcome variance that the TPE surrogate relies on to distinguish the effect of a hyperparameter choice from run-to-run noise. Reproducibility of data-loader shuffling within a trial is instead obtained by seeding the shuffling generator with the trial index.

\paragraph{Final model}
For each forecast horizon, the best-performing hyperparameter configuration found by the search above is retrained from scratch on the same training split, under the same early-stopping regime, and evaluated once on the held-out test split (Section~\ref{sec:evaluation}). The retraining is repeated for five random seeds (identical data and hyperparameters, independent weight initialisation and shuffle order) and results are reported as mean $\pm$ standard deviation across seeds, so that the reported spread reflects training noise rather than tuning quality. Training was performed on [INSERT: hardware/GPU specification].

\subsection{Evaluation protocol}
\label{sec:evaluation}

\paragraph{Baselines}
The reference baseline for skill scoring is persistence: the last observed spectrum at the end of the input window, broadcast unchanged across all $\tau$ forecast steps. This baseline uses only information available at forecast time and involves no model parameters. A second, stronger baseline is a linear autoregressive model, fitted independently per frequency bin: each bin's own history is regressed on its previous $L$ values (ridge-regularised, $10^{-6}$, on the training split only) and rolled out recursively for $\tau$ steps, feeding each prediction back as the newest lag — mirroring the Transformer's own autoregressive rollout. Because it has no cross-frequency coupling and no nonlinearity, it isolates how much of the model's skill comes from genuinely nonlinear, cross-bin structure rather than from per-bin temporal extrapolation.

\paragraph{Core metrics}
For every configuration, the following are computed both per forecast step and on all steps/samples flattened together: RMSE; Skill Score $SS = 1 - \mathrm{RMSE}_{\text{model}}/\mathrm{RMSE}_{\text{persistence}}$ (positive values indicate the model outperforms persistence; $SS=1$ is a perfect forecast); Pearson correlation coefficient; bias (mean signed
error); and $R^2$. For the spectral branches, all of these are frequency-weighted by the normalised trapezoidal weights introduced with Eq.~\eqref{eq:loss}, and are computed in the same log space in which the model is trained; physical-space counterparts of the per-step RMSE, skill-score and bias curves are computed separately (after exponentiation) whenever a comparison is made against a baseline or a model that operates in physical units.

Model selection — both across trials and across epochs within a trial — uses the Skill Score at the \emph{final} forecast step rather than an average across steps, since the intermediate autoregressive steps are scaffolding required to reach the chosen horizon rather than deliverables in their own right; an exponentially-weighted mean per-step Skill Score (weights halving at the midpoint of the horizon) remains implemented and is reported for comparability with the earlier configuration of this work, in which it was the selection criterion. For the composite-loss configuration of Eq.~\eqref{eq:loss}, however, selection uses a peak-fidelity criterion instead,
\begin{equation}
\text{PF} \;=\; \overline{\text{recall}} \;-\; \overline{\text{rel.\,err.}},
\end{equation}
the mean over the wind-sea and swell partition classes of the peak-separation recall minus the mean of the peak-height relative error (both defined below). The reason is structural: every Skill Score variant is a monotone transform of RMSE, and once $\lambda_{\text{base}} = 0$ the model is not trained on RMSE at all, so selecting by an RMSE-derived criterion would systematically favour whichever candidate drifted least from RMSE-friendly (i.e. blurred, hedged) behaviour, which is precisely the failure mode the composite loss exists to correct.

RMSE is not directly comparable across Optuna trials, because the input window length $L$ is itself tuned per trial and therefore the persistence baseline (and hence what counts as an "achievable" RMSE) differs from trial to trial, whereas Skill Score normalises for this. Notably, within a single trial (with a fixed $L$) SS and RMSE are equivalent since the persistence baseline is constant for that trial.

\paragraph{Bulk-parameter and spectral-shape diagnostics}
In addition to the core metrics above, the following are computed on the final forecast step after exponentiating predictions and targets back to physical units (and, where the physical spectrum is required, after denormalising $E = \tilde{E}\cdot\mu(f)$), with all spectral-moment integrals evaluated by the trapezoidal rule on the true,
non-uniformly spaced frequency grid (Section~\ref{sec:data}) rather than assuming a constant $\Delta f$:
\begin{itemize}
  \item RMSE, bias, and MAPE of $H_s = 4\sqrt{m_0}$ (predicted vs.\ target), obtained from the magnitude branch or, for the monolithic configuration, by integrating the forecast spectrum;
  \item RMSE and bias of $T_{m02} = \sqrt{m_0/m_2}$. For the shape branch this is computed directly on the unit-area spectrum: $T_{m02}$ is a ratio of moments and therefore invariant to a uniform rescaling of $E(f)$, so — unlike $H_s$ — it is physically meaningful without any magnitude information;
  \item spectral shape RMSE: both predicted and target spectra are normalised by the \emph{target} $m_0$ before
    computing per-spectrum RMSE, decoupling shape (frequency-distribution) error from energy-magnitude error, together with its Skill Score against a last-observed-shape persistence baseline. For the monolithic configuration, spectrum/step pairs whose target $m_0$ falls below $10^{-4}$~m$^2$ (corresponding to $H_s \approx 4$~cm) are excluded from this metric to avoid dividing by a near-zero denominator; this threshold was chosen to exclude only genuinely calm or missing records, well below the buoy's observed minimum $H_s$ ($\approx$0.8~m). The shape branch needs no such mask, since its targets are unit-area by construction;
  \item the mass-normalised $W_2$ distance of Eq.~\eqref{eq:loss}, reported as a metric alongside shape RMSE: it answers the same question with a shift-tolerant rather than pointwise geometry, so that "the peak is in the wrong place" and "the peak is missing" can be told apart by reading the two together. Its per-bin decomposition is likewise reported against that of RMSE;
  \item a mass-conservation check, the mean absolute deviation of $\int S_{\text{pred}}(f)\,\mathrm{d}f$ from unity across all steps, which should be zero by construction given the renormalisation in the rollout (Section~\ref{sec:architecture}) and is monitored as a regression guard;
  \item per-frequency-bin Scatter Index, $\mathrm{SI}(f) = \mathrm{RMSE}(f)/\overline{E_{\text{true}}(f)}$, computed over all samples at the final step, together with its frequency-weighted mean across bins, and the corresponding per-bin RMSE and bias curves.
\end{itemize}
$H_s$ MAPE is well-behaved for this record because the observed $H_s$ is
bounded away from zero ($\gtrsim$0.8~m), unlike a per-bin spectral MAPE,
which would be dominated by near-zero bins.

\paragraph{Peak-resolved, partition-conditioned diagnostics}
A frequency-weighted RMSE over the whole spectrum dilutes an error confined to a narrow peak (1--2 of 47 bins) almost to invisibility: a forecast that merges a distinct swell and wind-sea pair into one smoothed, lower hump can score as well as, or better than, a faithful one, because most of the weighted mass sits in the correctly predicted shoulders and tail. A dedicated set of peak-resolved metrics is therefore computed on the final forecast step.

Peaks are detected with the four-criterion significant-peak test of Portilla et al.~\cite{portilla2009}: a candidate local maximum is rejected if (i) its peak frequency exceeds $f_{\max} = 0.4$~Hz (high-frequency tail noise), (ii) the energy of its trough-to-trough partition is below a fraction $0.05$ of the spectrum's total energy, (iii) it has fewer than 2 spectral bins on either side before the next trough, or (iv) it is "sandwiched" between two higher-energy neighbouring peaks. This replaced an earlier scale-free prominence heuristic (a fraction of each spectrum's own maximum, tuned empirically) with published, physically motivated criteria. Each surviving partition is then classified as wind sea or swell by $\gamma^{*} = S_{\text{obs}}(f_p)/S_{\text{PM}}(f_p)$, the observed energy at the peak relative to the Pierson--Moskowitz reference spectrum evaluated analytically at that frequency, with $\gamma^{*} > 1$ denoting wind sea~\cite{violantecarvalho2009}.

All partition geometry and labels are taken from the \emph{true} spectrum only and applied to both sides. This mirrors the convention of the peak loss in Eq.~\eqref{eq:loss} and avoids partition-matching ambiguity: a partition the model misses entirely is simply absent from its peak list and is recorded as a recall miss, rather than being reconciled against some invented predicted window. The metrics are:
\begin{itemize}
  \item mean detected peak count, for true and predicted spectra separately — an over-smoothing model shows a systematically lower predicted count;
  \item peak-height relative error, $|\hat{E}(f_p) - E(f_p)|/E(f_p)$ evaluated at each true peak's bin. This is deliberately the harshest reading — it is not conditioned on the model having detected a peak there at all, so a fully blurred secondary peak still counts against it;
  \item peak-separation recall, the fraction of true peaks for which the model's own curve exhibits a significant local maximum (by the same criteria) within 2 bins — the direct measure of whether the model resolves distinct wave systems rather than merging them;
  \item $T_{m02}$ RMSE and bias computed \emph{within} each true partition's own trough-to-trough window, rather than over the whole spectrum. Partitions whose predicted in-window energy falls below the same $0.05$ fraction of their true in-window energy are excluded, since $T_{m02}$ is ill-conditioned on near-zero mass and a missed partition is already accounted for by recall.
\end{itemize}
Each of the last three is pooled twice: once overall, and once separately within each partition class. The class-conditional breakdown exists because pooling hides the attribution an ablation needs — wind-sea partitions are broad, energetic and fast-evolving (a magnitude-tracking problem), whereas swell partitions are narrow, slow and persistent (a position-tracking problem, precisely what the $W_2$ term targets) — so a change that helps only one of the two is invisible in the pooled number.

\paragraph{Correctness verification}
The normalise/denormalise round trip that underlies every physical-unit
metric above is verified independently of model training, using synthetic
JONSWAP spectra constructed so that $4\sqrt{m_0}$ equals a known target
$H_s$ by construction: the test suite confirms that (i) $H_s$ computed
directly from the physical spectrum recovers the target to within 2\%,
(ii) $H_s$ computed after a normalise/denormalise round trip through an
independently-generated $\mu(f)$ also recovers the target to within 2\%,
and (iii) $H_s$ computed from the \emph{normalised} spectrum directly
(i.e., omitting denormalisation) deviates from the target by more than 2\%,
confirming the denormalisation step is load-bearing rather than
incidental. The same suite verifies the unit-area property of the shape
transform, the analytic gradient identity of the soft peak-height term and
the zero-at-identity property of the KL and $W_2$ terms, the invariance of
the $W_2$ term to a uniform rescaling of either spectrum, the recovery of
known growth rates and oscillation frequencies by the DMD feature extractor
on synthetic data, and the boundary behaviour of the frequency-axis
convolution under replicate versus zero padding.

\subsection{Incremental ablation strategy}
\label{sec:ablation}

None of the components described above is taken as given. Each was introduced as a candidate, one at a time, and retained only if it improved the evaluation panel of Section~\ref{sec:evaluation} without degrading another axis of it, and only if the improvement justified its computational cost. This subsection states the protocol, the cost accounting that accompanies it, and what the protocol does and does not establish.

\paragraph{One change per iteration}
Each accepted change defines a new study version, recorded together with the input configuration, objective and architecture that produced it, so that results are never silently compared across incompatible settings. A change that alters the objective function, the target parameterisation or the search space invalidates comparison with earlier trials — the surrogate model of the sampler would otherwise be fitted to scores measured on two different scales — so such a change forces a fresh search under a new version, with the previous study database retained rather than overwritten. The sequence of accepted changes therefore forms an auditable chain, each link differing from its predecessor in one component: the frequency-structured embedding and the temporal convolutional front-end replacing a flat projection; the addition of $r_2$ as an input channel; the trapezoidal frequency weighting replacing a flat mean over the log-spaced grid; replicate padding at the frequency-axis boundary; attention pooling replacing the flatten-and-project aggregation; the switch to AdamW with a split dropout parameterisation; the logarithmic output parameterisation; the change of selection criterion from an exponentially-weighted mean to the final forecast step; the auxiliary dynamical features; and finally the composite loss.

\paragraph{Unit of comparison}
Because nearly every candidate interacts with capacity and regularisation — a richer embedding may want fewer layers, a sharper loss may want a different learning rate — a candidate is not transplanted into a frozen configuration. It is re-searched: the full hyperparameter search of Section~\ref{sec:hpo} is re-run with the candidate enabled, and the comparison is between the best configuration achievable with it and the best achievable without it. The cost of that choice is that a difference between two versions confounds the candidate with search noise. Two controls bound that confound. First, the selected configuration is retrained across five seeds (Section~\ref{sec:training}), giving a spread attributable to initialisation and shuffle order alone. Second, in controlled ablations the unchanged baseline arm is itself re-run five times, differing only in data shuffling, so that no arm is declared better until it separates from that spread.

\paragraph{Controlled ablation where a re-search would confound}
When the candidate space is small but the interaction risk is high, the protocol is inverted: the architecture is pinned and only the candidate is varied. This was the case for the loss terms of Eq.~\eqref{eq:loss}. Seven arms were run at a single horizon with every architectural and training hyperparameter pinned to one reference configuration: the per-bin baseline; $D_{\text{KL}}$, $W_2$ and $\mathcal{L}_{\text{peak}}$ each substituting the per-bin term alone; $W_2$ and $\mathcal{L}_{\text{peak}}$ each in combination with $D_{\text{KL}}$; and the three-way combination. Only the term weights were searched, one or two dimensions per arm. Re-searching the full space per arm was deliberately avoided: a budget of order ten trials cannot resolve a ten-dimensional space, and any difference recovered would confound the loss with a luckier draw of architecture.

The outcome illustrates why a single aggregate number is insufficient as an acceptance criterion. Relative to the per-bin baseline, the three-way combination improved peak-separation recall from $0.66$ to $0.94$, peak-height relative error from $0.41$ to $0.38$, partition-resolved $T_{m02}$ RMSE from $0.49$ to $0.41$~s and $T_{m02}$ bias from $-0.16$ to $+0.02$~s (each averaged over the two partition classes; seven of the eight class-resolved entries improved, the exception being an already near-zero wind-sea $T_{m02}$ bias). Every other arm traded at least one axis away. Most instructively, the peak term alone produced the highest recall of any arm ($0.97$) together with the \emph{worst} peak-height error ($0.65$ against a baseline of $0.41$) and the worst $T_{m02}$ error ($0.79$~s): optimising peak sharpness in isolation produces sharp peaks in the wrong places. An acceptance rule based on recall alone, or on any single scalar, would have adopted it.

\paragraph{Cost accounting}
Candidates are also judged on what they cost to compute, measured rather than assumed. Three regimes are distinguished.

The first is components whose cost is negligible against the forward and backward pass they accompany. For a representative batch of $128$ samples $\times$ $12$ forecast steps $\times$ $47$ bins, the per-bin weighted MSE, the KL term and the $W_2$ term cost approximately $0.3$, $0.5$ and $3.2$~ms per batch respectively (forward and backward, CPU), against a batch's own cost of the order of hundreds of milliseconds. All three are pure tensor operations with no Python-level loop and no optimal-transport solver, and $W_2$'s quantile inversion adds only a batched binary search and gather. The trapezoidal frequency weighting and the logarithmic target transform are likewise free at runtime. Components in this regime were adopted unconditionally once they showed an improvement, since keeping them costs nothing even where their benefit is small.

The second is components paid once rather than per step. The auxiliary dynamical features (Section~\ref{sec:data}) require one eigendecomposition per sample, measured at $185\,\mu$s per sample, and are computed during data preparation rather than during training: for the training split this is a fixed cost of a few seconds, amortised over up to 100 epochs, i.e.\ negligible. The per-frequency normalisation statistics are computed once on the same schedule. That a feature can be moved into this regime is itself a reason to prefer it: an equally informative feature that had to be recomputed per batch would have to clear a much higher bar.

The third is components that are genuinely expensive, and these are gated rather than enabled by default. The soft peak-height term requires trough-to-trough windows obtained from a non-differentiable peak detector that loops in Python over every spectrum: approximately $0.1$~ms per spectrum, or $\sim$$150$~ms per batch — about two orders of magnitude more than the three differentiable loss terms combined, and roughly $14$~s per epoch at the reference configuration, i.e.\ tens of minutes over a full trial. The effect is visible in the ablation's wall-clock: arms containing the peak term ran at roughly $30$--$36$ minutes per trial against roughly $20$ for the baseline arm, although that figure also reflects differences in how early trials were pruned. This term was therefore retained only because the ablation showed that the recall gain is not obtainable without it, and it is gated behind a non-zero weight, so that any configuration not using it pays nothing. The same gating applies on the evaluation side: the peak-resolved metrics are opt-in and are computed on the final forecast step only, reducing their cost to a fraction of a second per validation pass rather than the per-batch, per-step cost the training term incurs.

A measured cost may also constrain the search space instead of the model. The scheduled-sampling rollout unrolls one decoder pass per forecast step and backpropagates once, so its memory footprint grows quadratically with the horizon; rather than abandoning scheduled sampling at long horizons, the batch-size range is restricted for $\tau > 24$~h. Similarly, the encoder output is computed once per sample and reused across all $\tau$ decode steps, which makes autoregressive rollout affordable at evaluation time without changing the model. Pruning serves the same purpose at the level of the search: its warmup and patience settings (Section~\ref{sec:hpo}) were themselves tuned after a majority of trials at one horizon were found to be terminated at exactly the boundary epoch, wasting budget on trials that would have recovered.

\paragraph{Candidates tested and not retained}
Several candidates did not survive this procedure, and are reported because a description of the retained components alone would overstate how directed the path was. A spectral-slope (derivative) penalty, aimed at the same multimodal blurring problem as $W_2$, was implemented and reverted. Enforcing non-negativity of the spectrum architecturally, with a Softplus output activation during training and clipping at inference, was superseded by the logarithmic output parameterisation, which obtains the same guarantee by construction and without a saturating activation. Peak detection was initially based on a scale-free prominence heuristic — a fraction of each spectrum's own maximum, with the fraction tuned empirically on one station's test split — and was replaced by the published four-criterion test of Section~\ref{sec:evaluation}, which introduces no dataset-tuned constant. The wind auxiliary channel is implemented and was trialled, but cannot be assessed at the station reported here: every wind-direction and wind-speed entry in that station's standard meteorological record is a missing-value sentinel, so the channel is unavailable for this record after sentinel masking, and the dynamical features described above — derived from the spectra themselves and requiring no external data stream — are used in its place. Finally, the attention head dimension and head count were withdrawn from the search space rather than rejected: a tally across previous searches found the same two values winning at six of eight (study, horizon) points, whereas every other hyperparameter spanned nearly its full range with no consensus, so pinning them reallocates budget to the dimensions that remain undetermined.

\paragraph{Limitations of the protocol}
Two caveats follow directly from the design. Because each accepted change is re-searched rather than transplanted, attribution of a gain to a single component is only as strong as the noise estimate that accompanies it; differences of the order of the seed-to-seed spread are reported as such and not claimed as improvements. And because the peak-resolved, partition-conditioned panel was introduced partway through the sequence, components adopted before it were accepted on the evidence of RMSE-family metrics alone and have not been re-examined against the panel — which, as the loss ablation demonstrates, is exactly the kind of evidence that can favour an over-smoothed forecast. Re-testing the earlier architectural decisions under the current panel remains outstanding.
